\section{Conclusion}\label{sec:conclusion}

In this work, we presented \sys, a comprehensive, dynamic, safe, and efficient extension framework. \sys is based upon two design principles: maintaining compatibility with the existing eBPF ecosystem and executing userspace extensions in the extended application without requiring kernel involvement.  We show that \sys full system extensions are useful across \numcases case studies and that the system provides a large performance improvement compared to existing state-of-the-art full system extension tools.  We  will release \sys as an open-source project. 

%, a cutting-edge userspace eBPF runtime designed to alleviate the identified performance and security drawbacks inherent to kernel eBPF. Through \textit{bpftime}, we significantly improved Uprobe and Syscall hook performance, demonstrating a tenfold reduction in Uprobe overhead relative to kernel uprobes. Our runtime showcases a smooth integration with ongoing processes, bypassing the need for restarts or manual recompilations. Compatibility with prevalent eBPF toolchains like clang and libbpf simplifies userspace eBPF development, portraying \textit{bpftime} as a step forward in enhancing userspace eBPF runtimes. The open-sourced nature of \textit{bpftime} at \href{https://github.com/eunomia-bpf/bpftime}{https://github.com/eunomia-bpf/bpftime} paves the way for community engagement, fostering further exploration and innovation in this domain.